\section{Introduction}
\label{sec:intro}

Application-level autonomous agent systems such as OpenClaw~\citep{openclaw-github} have matured from research demonstrations into real-world workers, deployed across Slack, Discord, web, and other channels to complete complex multi-step tasks for users. Yet once deployed, they remain largely static: they do not learn from how they are actually used, and the same failure modes recur across users until the next human-driven iteration ships a fix. A wave of self-evolving agents has emerged in response, letting agents evolve themselves after deployment~\citep{hermes-github, hermes-self-evo, wang2026procedural, ma2026skillclaw, liang2026genericagent, wang2025evoagentx}.

These systems leave one layer untouched. As Table~\ref{tab:scope-comparison} makes visible, their evolution scope is bounded to \textbf{text-mutable artifacts}---skill files, prompt configurations, memory schemas, and at most workflow graphs; the \textbf{agent harness}---routing, state management, dispatch, hooks, mediator, session lifecycle---is never modified by the agent itself. The limit this boundary imposes is physical: edits to text-mutable artifacts can only change \emph{what the agent itself thinks} (what to say, which skill to invoke, how to break a task into sub-goals); they cannot change \emph{what the harness decides on its behalf}. Once a failure originates in this layer---mis-routed messages, hooks firing out of order, corrupted session state, atomicity bugs across concurrent skills---no update to skills, prompts, or memory can reach it: the bug is not in the prompt text, and a prompt rewrite cannot paper over it. The proportion of failures of this kind scales with harness complexity, so this gap widens as agentic systems mature.

\begin{table}[ht]
\caption{Scope of evolution across application-level self-evolving agentic systems.}
\label{tab:scope-comparison}
\centering
\small
\begin{tabular}{lcccc}
\toprule
Project & Skill & Prompt & Memory & Harness \\
\midrule
Hermes Agent~\citep{hermes-github}     & $\checkmark$ & $\times$              & $\checkmark$ & $\times$ \\
SkillClaw~\citep{ma2026skillclaw}            & $\checkmark$ & $\times$              & $\times$     & $\times$ \\
GenericAgent~\citep{liang2026genericagent}             & $\checkmark$ & $\times$              & $\checkmark$ & $\times$ \\
EvoAgentX~\citep{wang2025evoagentx}                 & $\checkmark$ & $\checkmark$          & $\times$     & $\times$ \\
\textbf{MOSS (Ours)}                             & \boldmath$\checkmark$ & \boldmath$\checkmark$ & \boldmath$\checkmark$ & \boldmath$\checkmark$ \\
\bottomrule
\end{tabular}
\end{table}

We argue that source-level adaptation---where an agent performs self-rewriting on its own source---is a fundamentally more general medium, simultaneously superior to text-mutable evolution along four axes. It is Turing-complete: because programming languages are Turing-complete, the source-code design space constitutes a universal search space within which every text-mutable agent-design space---prompts, skills, memory schemas, workflow graphs---sits as a strict subset~\citep{hu2025automated}; editing code can reach not only every configuration these subsets can represent, but arbitrary agentic structures they cannot. It is a strict superset of every text-mutable scope: whatever a prompt edit can achieve, an equivalent code edit can also achieve, and not the other way around. It takes effect deterministically: routing logic, hook ordering, and state-machine invariants run as code, and their behavior does not depend on whether the base model correctly reads new text and complies with it---text-mutable fixes are the opposite, with effectiveness rising and falling with current base-model capability. And it does not erode under long-context drift: text-mutable fixes are prompts, skills, and memory entries that the agent must re-read on every turn, and as these accumulate over weeks of production use the model's adherence to any single piece of guidance dilutes; edits at the source layer are encoded as behavior, not text to be re-read, and so do not degrade as the system ages.

Academic work has shown source-level self-rewriting to be feasible on minimal scaffolds: SICA~\citep{robeyns2025self}, the Darwin G\"{o}del Machine~\citep{zhang2025darwin}, and HyperAgents~\citep{zhang2026hyperagents} all demonstrate an agent modifying its own implementation to lift its own benchmark scores; Meta-Harness~\citep{lee2026meta} further shows that exposing past execution traces to the coding-agent proposer drives iterative improvement more strongly than benchmark scores alone do. But these systems all run on minimal scaffolds with feedback gated by benchmark scores---fundamentally an exploratory paradigm. Carrying the same source-level capability over to an application-level, production-grade substrate is a different engineering setting altogether: the codebase is large, there is no clean benchmark score to anchor evolution against, and live user traffic together with persistent user state must survive. This pushes a series of concrete questions up to the system-design layer: when to trigger an evolution attempt, where in the codebase the fix should land, whether the resulting candidate is actually better than the deployed version, and how to put that candidate in front of real users without disrupting the live deployment.

We present \textbf{MOSS}, a system that performs comprehensive source-level self-rewriting on production agentic substrates---an instantiation of the source-level adaptation argued for above. As Table~\ref{tab:scope-comparison} shows, it is the only system that reaches the harness, covering a strict superset of the text-mutable scope addressed by prior work. MOSS surfaces its entire evolution lifecycle---triggering, status query, stop, apply, conversational flagging---to the agentic substrate as a built-in \texttt{moss evo} CLI capability via system-prompt injection, so the user interacts with evolution through the same conversational interface they already use for daily work. Evidence accumulates from two complementary channels: a periodic background scan over recent user sessions automatically surfaces under-performing dialogue segments, and the user can additionally flag any turn in conversation; both feed the same batch, and every downstream stage is anchored to fixing that batch rather than to a synthetic benchmark target. Code modification is carried out by a pluggable external coding-agent CLI invoked as a host-side subprocess, and candidates are verified on ephemeral trial workers that replay the batch against the candidate image in a production-equivalent container---not against unit tests or synthetic harnesses. Throughout the loop, every artifact---plans, diffs, build logs, per-iteration scoring matrices---is readable by the agent on demand, so the user can audit what the system intends to change and why before authorizing the swap. On convergence, MOSS proactively notifies the user through the same conversational channel and pauses; only after explicit user authorization does the host-daemon perform an in-place container swap, preserving user state and rolling back automatically on a failed post-swap health probe. On OpenClaw, MOSS lifts a four-task mean grader score from 0.25 to 0.61 in a single cycle, without human intervention.

